\chapter{Interface utilisateur (Blazor)}
\label{ch:ui}

\section{Structure de navigation}

L'application propose un menu latéral avec les entrées : tableau de bord, actifs, transactions, budget, outils (marchés, import CSV, backtest). Le layout principal inclut un commutateur de \textbf{thème clair/sombre} (attribut \texttt{data-bs-theme} et stockage \texttt{localStorage}).

\section{Pages principales}

\begin{description}
  \item[Tableau de bord] Cartes KPI (valeur totale, liquidité, gain/perte, ROI), alertes, graphiques Chart.js (répartition en anneau, courbe d'historique de valeur), répartition par actif et par type, bloc analytique par symbole, tableau de performance avec CMP et coût de revient.
  \item[Actifs] Formulaire CRUD, filtres, pagination, tri, recherche avec \emph{debounce}.
  \item[Transactions] Saisie avec validation, résumé budget, historique paginé, tri, export CSV global.
  \item[Budget] Ajustement du budget initial et de la liquidité.
  \item[Outils] Documentation des modes de prix, bouton de rafraîchissement API, zone de texte pour import CSV, formulaire de backtest.
\end{description}

\section{Technologies front}

\begin{itemize}[leftmargin=*]
  \item \textbf{Bootstrap} (fichiers locaux sous \texttt{wwwroot/lib/bootstrap}) pour la grille et les composants.
  \item \textbf{Chart.js} chargé via CDN dans \texttt{App.razor}, scripts personnalisés \texttt{portfolio-charts.js} et \texttt{theme.js}.
  \item \textbf{Interop JavaScript} (\texttt{IJSRuntime}) depuis les composants interactifs pour initialiser et mettre à jour les graphiques.
\end{itemize}

\section{Mode de rendu}

Les pages à interactions fortes utilisent \texttt{@rendermode InteractiveServer}. Le reconnect modal Blazor est conservé pour une meilleure expérience en cas de perte de circuit SignalR.
